\documentclass[main.tex]{subfiles}


\begin{document}

%from pagenumber to up
% \phantomsection 
% \hypertarget{toc}{}

 

 

 
   

\section{Фотоэффект}


Эйнштейн в поисках объяснения результатов опытов по взаимодействию света\footnote{Здесь и далее любое электромагнитное излучение называется для краткости светом.} с веществом пришел к выводу, что \emph{свет (и электромагнитное излучение вообще) состоит из отдельных порций~--- квантов} (или, как говорят, \emph{фотонов})\footnote{Эйнштейн развил идею Планка (\emph{гипотезу Планка}), которая состоит в том, что свет излучается и поглощается отдельными порциями (квантами).}. Таким образом, стали считать, что свет~--- это поток особых частиц (фотонов), движущихся в вакууме со скоростью~$c$.

\textbf{Энергия фотона} пропорциональна частоте излучения:
\begin{equation}
    E_\text{ф}=h\nu,
\end{equation}

\sbox{0}{%
    \begin{tikzpicture}[font=\small,            line cap=round,                             >={Stealth[inset=0pt,round,length=3mm, angle'=20]},vec/.style={>={Stealth[inset=0pt,round,length=3.5mm, angle'=20]}}]


\tikzset{em/.pic={
  \begin{scope}[shift={({rnd*360}:.025)}]
  \begin{scope}[transparency group,nearly transparent]
  \begin{scope}[transparency group,semitransparent]
  \shade[ball color=NavyBlue!70] ([shift={(0:\de cm)}]0,0) circle (\de);
  \end{scope}
  \begin{scope}[transparency group]
  \shade[ball color=NavyBlue!70] ([shift={(180:{0 cm})}]0,0) circle (\de);
  \end{scope}
  \end{scope}
  \shade[ball color=NavyBlue!70] ([shift={(180:{\de cm})}]0,0) circle (\de);
  \end{scope}}
}

\tikzset{ph/.pic={
  \begin{scope}[shift={({rnd*360}:.08)}]
  \begin{scope}[transparency group,nearly transparent]
  \begin{scope}[transparency group,semitransparent]
  \shade[ball color=yellow!60!orange] ([shift={({90+\an}:\de cm)}]0,0) circle (\de);
  \end{scope}
  \begin{scope}[transparency group]
  \shade[ball color=yellow!60!orange] ([shift={(180:{0 cm})}]0,0) circle (\de);
  \end{scope}
  \end{scope}
  \shade[ball color=yellow!60!orange] ([shift={({\an-90}:{\de cm})}]0,0) circle (\de);
  \end{scope}}
}

\def\de{.1}
\def\an{40}


% plates
\fill[red,semitransparent] (0,0) rectangle++ (.1,2);
\draw[] (0,0) rectangle++ (.1,2);
\draw (0.05,0) node[yshift=2 cm,above] {А} to[*short,-o] ++ (0,-1.5) node[below] {$+$} coordinate (ll);
\draw[->,vec,thick,red] (0.05,0) ++ (0,-.75) ++ (-.2,0) ++ (0,-.5) --++ (0,1) node[midway,black,left] {$I_\text{ф}$}; 

\fill[NavyBlue,semitransparent] (3,0) rectangle++ (.1,2);
\draw[] (3,0) rectangle++ (.1,2);
\draw (3.05,0) node[yshift=2 cm,above] {К} to[*short,-o] ++ (0,-1.5) node[below] {$-$} coordinate (rr);
\draw[<-,gray,shorten <=1.5mm,shorten >=3mm] (ll) --++ (1.5,0) node[black] {$U$};
\draw[<-,gray,shorten <=1.5mm,shorten >=3mm] (rr) --++ (-1.5,0);


% light
\fill[yellow,semitransparent] (3,.4) --++ ({90+\an}:3) --++ ({\an}:{1.2*cos(\an)}) -- (3,1.6) -- cycle;
\pic at (1.8,2.7) {ph};
\pic at (1.9,2) {ph};
\pic at (2.7,1.5) {ph};


% e
\pic at (2,1.5) {em};
\pic at (2.7,1) {em};
\pic at (2.4,.5) {em};





    \end{tikzpicture}%
}

{%
\setlength\intextsep{0pt}
\begin{wrapfigure}[]{r}{\wd0}
    \centering
    \usebox{0}%
    \caption{Фотоэффект}
    \label{fig:p197}
\end{wrapfigure}

\noindent где $h$~--- \emph{постоянная Планка} (см.~справочные таблицы).


%     \begin{figure}[H]
%     \centering

%     \begin{tikzpicture}[font=\small,            line cap=round,                             >={Stealth[inset=0pt,round,length=3mm, angle'=20]},vec/.style={>={Stealth[inset=0pt,round,length=3.5mm, angle'=20]}}]


% \tikzset{em/.pic={
%   \begin{scope}[shift={({rnd*360}:.025)}]
%   \begin{scope}[transparency group,nearly transparent]
%   \begin{scope}[transparency group,semitransparent]
%   \shade[ball color=NavyBlue!70] ([shift={(0:\de cm)}]0,0) circle (\de);
%   \end{scope}
%   \begin{scope}[transparency group]
%   \shade[ball color=NavyBlue!70] ([shift={(180:{0 cm})}]0,0) circle (\de);
%   \end{scope}
%   \end{scope}
%   \shade[ball color=NavyBlue!70] ([shift={(180:{\de cm})}]0,0) circle (\de);
%   \end{scope}}
% }

% \tikzset{ph/.pic={
%   \begin{scope}[shift={({rnd*360}:.08)}]
%   \begin{scope}[transparency group,nearly transparent]
%   \begin{scope}[transparency group,semitransparent]
%   \shade[ball color=yellow!60!orange] ([shift={({90+\an}:\de cm)}]0,0) circle (\de);
%   \end{scope}
%   \begin{scope}[transparency group]
%   \shade[ball color=yellow!60!orange] ([shift={(180:{0 cm})}]0,0) circle (\de);
%   \end{scope}
%   \end{scope}
%   \shade[ball color=yellow!60!orange] ([shift={({\an-90}:{\de cm})}]0,0) circle (\de);
%   \end{scope}}
% }

% \def\de{.1}
% \def\an{40}


% % plates
% \fill[red,semitransparent] (0,0) rectangle++ (.1,2);
% \draw[] (0,0) rectangle++ (.1,2);
% \draw (0.05,0) node[yshift=2 cm,above] {А} to[*short,-o] ++ (0,-1.5) node[below] {$+$} coordinate (ll);
% \draw[->,vec,thick,red] (0.05,0) ++ (0,-.75) ++ (-.2,0) ++ (0,-.5) --++ (0,1) node[midway,black,left] {$I_\text{ф}$}; 

% \fill[NavyBlue,semitransparent] (3,0) rectangle++ (.1,2);
% \draw[] (3,0) rectangle++ (.1,2);
% \draw (3.05,0) node[yshift=2 cm,above] {К} to[*short,-o] ++ (0,-1.5) node[below] {$-$} coordinate (rr);
% \draw[<-,gray,shorten <=1.5mm,shorten >=3mm] (ll) --++ (1.5,0) node[black] {$U$};
% \draw[<-,gray,shorten <=1.5mm,shorten >=3mm] (rr) --++ (-1.5,0);


% % light
% \fill[yellow,semitransparent] (3,.4) --++ ({90+\an}:3) --++ ({\an}:{1.2*cos(\an)}) -- (3,1.6) -- cycle;
% \pic at (1.8,2.7) {ph};
% \pic at (1.9,2) {ph};
% \pic at (2.7,1.5) {ph};


% % e
% \pic at (2,1.5) {em};
% \pic at (2.7,1) {em};
% \pic at (2.4,.5) {em};





%     \end{tikzpicture}
% \caption{Фотоэффект}
% \label{fig:p197}
% \end{figure}


\textbf{Фотоэффект}~--- это вырывание электронов из тела падающим светом. Явление фотоэффекта было иссследовано Столетовым с помощью специального устройства, основная часть которого показана на рис.\,\ref{fig:p197}. 

Главная часть этого устройства состоит из двух металлических пластин~--- пластина~А (анод) и пластина~К (катод). К пластинам от батарейки подводится напряжение~$U$ (на анод подан <<плюс>>, а на катод~--- <<минус>>). В данном случае напряжение~$U$ считается положительным (его знак определяется знаком анода).

При освещении катода, например, ультрафиолетовым светом его фотоны (желтые шары) выбивают с катода электроны (синие шары), которые разгоняются электрическим полем (созданным пластинами) в сторону анода. Достигшие анода выбитые электроны (\emph{фотоэлектроны}) устремляются к <<плюсу>>~--- через батарейку протекает ток~$I_\text{ф}$, называемый \emph{фототоком} (потому что его создают фотоэлектроны). 

% происходит перенос отрицательного заряда с пластины~К на пластину~А, что эквивалентно протеканию тока

}

Три \emph{закона фотоэффекта} формулируются так.
\begin{enumerate}
    \item[\textbf{I.}] Число электронов, вырываемых из катода за секунду, пропорционально мощности падающего на катод излучения (при его неизменной частоте)\footnote{Мощность падающего излучения~--- это отношение суммарной энергии фотонов, попадающих на облучаемое тело, ко времени, за которое эти фотоны попали на тело.}:
    \begin{equation}
        N_\text{фэ}\sim P_\text{изл}.
    \end{equation}
    \item[\textbf{II.}] Максимальная кинетическая энергия фотоэлектронов линейно возрастает с частотой света и не зависит от мощности падающего излучения\footnote{Эту зависимость можно использовать только для качественных ответов: например, чем больше частота излучения, тем больше максимальная кинетическая энергия фотоэлектронов.}:
    \begin{equation}
        E_\text{к.\,max}\sim \nu.
    \end{equation}
    \item [\textbf{III.}] Фотоэффект наблюдается только при частотах, больших некоторой частоты~$\nu_\text{кр}$, называемой \emph{красной границей фотоэффекта} и зависящей от облучаемого вещества:
    \begin{equation}
        \nu>\nu_\text{кр}.
    \end{equation}
    
    % Для каждого вещества существует \emph{красная граница фотоэффекта}~--- наименьшая частота света~$\nu_\text{кр}$, при которой фотоэффект ещё возможен. При $\nu < \nu_\text{кр}$ фотоэффект не наблюдается ни при какой мощности падающего света.
   
\end{enumerate}








 
\end{document}